\pg{355} outro, exige três elementos: --- 1º.) um \emph{principio}, de que se parte e que importa alguma coisa, a qual se perde; --- 2º.) um \emph{termo}, a que se chega e que importa alguma coisa, a qual se adquire; --- 3º.) um \emph{sujeito}, que do \emph{principio} chega ao \emph{termo} e se encontra num e noutro extremo. --- Se faltar algum d' estes elementos, não haverá verdadeira \emph{mudança} \footnote{O \emph{principio}, de que se parte, é chamado, pelos escholasticos, termo a \emph{quo}; o \emph{termo}, a que se chega, termo \emph{ad} quem. --- Entre o \emph{principio} e o \emph{termo}, ou entre o termo a \emph{quo} e o termo \emph{ad} quem, deve existir uma certa opposição; aliás não poderia conceber-se a mudança. Diz S. Thomaz: « \emph{Omnis mutatio est inter opposita} » (\emph{In V Physic.}, l. 2). --- A mudança recebe a espécie e a denominação --- não do \emph{principio}, ou do termo a \emph{quo}, --- mas do \emph{termo}, ou do termo \emph{ad quem}, porque é para este termo que a mudança tende e é nelle que se completa; assim dizemos que a água está quente, quando de fria se muda para quente. Escreve o S. Doutor: « \emph{Mutationes accipiunt speciem et dignitatem non a termino a quo, sed a termino ad quem} » (\emph{Sum. Th.}, p. I, q. 45, a. 1 ad 2). --- O \emph{sujeito} da mudança deve ser um e o mesmo no \emph{principio} e no \emph{termo}; o que muda propriamente é o \emph{estado} do sujeito. Esta mudança do estado realisa-se emquanto o sujeito perde uma coisa e adquire outra. --- Se, pois, o sujeito não se encontrasse no \emph{principio}, ou no \emph{termo}, não haveria propriamente \emph{mudança}. Assim não é mudança --- nem a \emph{creação}, porque antes d' ella não existe o sujeito, isto é, não há um sujeito commum aos dois extremos, --- nem a \emph{annihilação}, porque depois d' ella não existe o sujeito, isto é, não há um sujeito commum aos dois extremos, pois o \emph{nada} não é um \emph{estado}, mas uma mera negação (\emph{Cf. C. Gent.}, II, 17; \emph{De Pot.}, q. 3, a. 2).}.

\textbf{§4. Espécies de mudança.} --- A \emph{mudança} é --- \emph{intrinseca} ou \emph{extrinseca}, --- \emph{instantanea} ou \emph{successiva}, --- \emph{substancial} ou \emph{accidental}.

a) \emph{Mudança intrinseca} ou \emph{extrinseca}. --- É \emph{intrinseca}, quando o sujeito adquire ou perde alguma coisa em si mesmo; tal é a mudança de um corpo, que de frio se torna quente. --- É \emph{extrinseca}, quando o sujeito nada adquire ou perde em si mesmo, mas apenas reveste uma nova relação por se ter mudado um outro ente, que de algum modo lhe se refere; assim dizemos que uma torre, que estava á nossa direita, se muda emquanto passa para a nossa esquerda pelo facto de termos nós mudado de posição \footnote{Esta distinção não deve confundir-se com a distinção da mudança proveniente de uma causa \emph{intrinseca}, ou de uma causa \emph{extrinseca}. --- Mudança proveniente de uma causa \emph{intrinseca} é a que se funda na própria natureza de sujeito, que tem a potencialidade para a não-existencia, como é toda a substancia composta de materia e de forma. --- Mudança proveniente de uma causa \emph{extrinseca} é a que tem por \emph{principio} a potencia activa de um agente exterior, de que o sujeito depende; assim a nossa alma, embora não tenha\pgn{356}em si a potencia para a não-existencia, comtudo, por estar sujeita ao poder omnipotente de Deus, pode mudar. --- Como se vê, a mudança \emph{intrinseca} pode derivar de uma causa \emph{extrinseca}.}
